\documentclass[12pt]{article}

\setlength{\parindent}{0pt}
\setcounter{secnumdepth}{3}
\usepackage[margin=1in]{geometry}

% packages
\usepackage{amsmath, amsfonts, amsthm, amssymb}
\usepackage{graphicx, float}
\usepackage{mathtools}
\usepackage{titlesec}
\usepackage{interval}
\usepackage{hyperref}
\usepackage{siunitx}
\usepackage{titling}
\usepackage{vwcol}
\usepackage{setspace}
\usepackage{empheq}
\usepackage{cancel}
\usepackage{esdiff}
\usepackage{multicol}
\usepackage{mdframed}
\usepackage{esdiff}
\usepackage{tikzsymbols}
\usepackage{multicol}
\usepackage{tikz}
\usepackage{varwidth}
\usepackage{pgfplots}
\usepackage{fancyhdr}
\usepackage{enumitem}
\usepackage{tcolorbox}

\intervalconfig {
	soft open fences
}

% header
\pagestyle{fancy}
\fancyhf{}
\lhead{Ethan Fan}
\rhead{AP Calculus BC}
\cfoot{\thepage}

% section format
\titleformat{\section}
  {\large \bfseries}
  {}
  {0em}
  {}
\titlespacing*{\section}{0pt}{0.5em}{0.3em}

\titleformat{\subsection}[runin]
  {\bfseries}
  {}
  {0em}
  {}
\titlespacing*{\subsection}{0pt}{0pt}{0em}

\titleformat{\subsubsection}[runin]
  {\itshape}
  {(\roman{subsubsection})}
  {0em}
  {}

% commands
\newcommand{\alignedintertext}[1]{%
  \noalign{%
    \vskip\belowdisplayshortskip
    \vtop{\hsize=\linewidth#1\par
    \expandafter}%
    \expandafter\prevdepth\the\prevdepth
  }%
}

\newcommand*{\paren}[1]{\ensuremath\left(#1\right)}
\newcommand*{\limit}[2][x]{\ensuremath{\displaystyle\lim_{#1\to#2}}}
\newcommand*{\Limit}[3][x]{\ensuremath{\displaystyle\lim_{#1\to#2}\left[#3\right]}}
\newcommand*{\deriv}[1][x]{\ensuremath{\dfrac{\mathrm{d}}{\mathrm{d}#1}}}
\newcommand*{\Deriv}[2][x]{\ensuremath{\dfrac{\mathrm{d}}{\mathrm{d}#1}\left[#2\right]}}
\newcommand{\ddx}{\dfrac{d}{dx}}
\newcommand{\dydx}{\dfrac{dy}{dx}}
\newcommand{\dydt}{\dfrac{dy}{dt}}
\newcommand{\dxdt}{\dfrac{dx}{dt}}
\newcommand*{\iinteg}[2][x]{\ensuremath{\displaystyle\int #2\;\mathrm{d}#1}}
\newcommand*{\dinteg}[4][x]{\ensuremath{\displaystyle\int_{#2}^{#3}#4\;\mathrm{d}#1}}
\newcommand*{\abs}[1]{\ensuremath{\left|#1\right|}}
\newcommand*{\eps}{\varepsilon}
\newcommand*{\real}{\ensuremath\mathbb{R}}
\newcommand*{\integer}{\ensuremath\mathbb{Z}}
\newcommand*{\posint}{\ensuremath\mathbb{Z^+}}
\newcommand*{\cbrt}[1]{\ensuremath{\sqrt[3]{#1}}}
\newcommand*{\lheqzero}{\ensuremath{\underset{\text{L'H}}{\overset{\left[\frac00\right]}{=}}}}
\newcommand*{\lheqinfty}{\ensuremath{\underset{\text{L'H}}{\overset{\left[\frac{\infty}{\infty}\right]}{=}}}}
\newcommand{\tor}{\text{ or }}
\newcommand{\floor}[1]{\lfloor #1 \rfloor}

\newcommand{\vem}{\vspace{0.5em}}
\newcommand{\example}[1]{\section{Example #1}}
\newcommand{\problem}[1]{\section{Problem #1}}
\newcommand{\subproblem}[1]{\subsection{(#1)} \,}

\DeclareMathOperator{\dne}{DNE}
\DeclareMathOperator{\sgn}{sgn}

\newcommand{\definition}[1]{\begin{tcolorbox}[colback=red!5!white,colframe=red!75!black,parbox=false] #1 \end{tcolorbox}}

\newcommand{\theorem}[2]{\begin{tcolorbox}[title={#1},colback=blue!5!white,colframe=blue!75!black,parbox=false] #2 \end{tcolorbox}}

\allowdisplaybreaks
% \postdisplaypenalty=100000

% document
\begin{document}

\vspace*{0cm}

% opening
\begin{center}
    {\Large \bfseries Problem Set 21} \\[0.5cm]
    {\large Volumes of Revolution II: Other Axes} \\[0.35cm]
    {\normalsize October 27, 2026}
\end{center}

% problems

\problem{9}

\subproblem{a}
\begin{gather*}
    S=\int_0^33x-x^2\, dx = \frac{3}{2}x^2-\frac{x^3}{3}\Big|^3_0=\boxed{\frac{9}{2}}
\end{gather*}

\subproblem{b}
\begin{gather*}
    V=\pi \int_0^3 (4x-x^2+2)^2-(x+2)^2\, dx\\
    =\pi \int_0^3 x^4+16x^2+4-4x^2+16x-8x^3-x^2-4x-4\, dx\\
    \pi (\frac{x^5}{5}-2x^4+\frac{11}{3}x^3+6x^2)\Big|_0^3\\
    = \boxed{\frac{198\pi}{5}}
\end{gather*}

\subproblem{c}
\begin{gather*}
    V = \pi \int_{0}^{3} \left(2 + \sqrt{4-y}\right)^2 - (4-y)^2  dy \\+ \pi \int_{3}^{4}  \left(2 + \sqrt{4-y}\right)^2 - \left(2 - \sqrt{4-y}\right)^2 dx
\end{gather*}

\subproblem{d}
\begin{gather*}
    W(k) = \pi \int_{0}^{3} (f(x) + k)^2 - (g(x) + k)^2 dx\\
    W'(k)=2\pi \int_0^3 f(x)-g(x)\, dx\\
    =\boxed{9\pi}
\end{gather*}

\subproblem{e}
\begin{gather*}
    V=2\pi (k-y)A\\
    \frac{198\pi}{5}=2\pi (y+2)\frac{9}{2}\\
    22=5y+10 \implies y= \frac{12}{5}\\
    k=\frac{12}{5}+\frac{22}{5}=\boxed{\frac{34}{5}}
\end{gather*}

\problem{10}

\subproblem{a}
\begin{gather*}
    T=\int_0^2\frac{4}{4+x^2} \, dx = 4\cdot \frac{1}{2} \arctan 1 = \boxed{\frac{\pi}{2}}
\end{gather*}

\subproblem{b}
\begin{gather*}
    T=\pi \int_0^2\paren{\frac{4}{4+x^2}}^2 \, dx\\
    =\pi \int_0^2 \frac{16}{(4+x^2)^2}\,dx\\
    = 16\pi \int_0^\frac{\pi}{4} \frac{2\sec^2\theta\, d\theta }{16\sec^4 \theta }\\
    = \pi \int_0^\frac{\pi}{4} 2\cos^2 \theta \,d\theta \\
    =\pi \int_0^\frac{\pi}{4} \cos 2\theta +1\,d\theta\\
    = \pi (\frac{\sin 2\theta}{2}+\theta )\Big|_0^\frac{\pi}{4}\\
    = \boxed{\frac{\pi^2}{4}+\frac{\pi}{2}}
\end{gather*}

\subproblem{c}
\begin{gather*}
    y=\frac{4}{4+x^2} \implies x=2\sqrt{\frac{1}{y}-1}\\
    V=\pi \int_0^\frac{1}{2}4\, dy+\pi \int_{\frac{1}{2}}^1 \frac{4}{y}-4\, dy\\
    =2\pi + \pi (4\ln y )\Big|^1_\frac{1}{2}-2\pi\\
    =\boxed{4\pi \ln 2}
\end{gather*}

\subproblem{d}
\begin{gather*}
    V=\pi \int_0^21-(1-\frac{4}{4+x^2})^2\, dx\\
    =\pi \int_0^2 2\cdot \frac{4}{4+x^2} - \paren{\frac{4}{4+x^2}}^2\, dx\\
    = 2\pi \cdot  \frac{\pi}{2} - \frac{\pi^2}{4}-\frac{\pi}{2}\\
    = \boxed{\frac{3\pi^2}{4}-\frac{\pi}{2}}
\end{gather*}

\subproblem{e}
\begin{gather*}
    T=\int_0^b\frac{4}{4+x^2} \, dx = 4\cdot \frac{1}{2} \arctan \frac{b}{2} = \frac{\pi}{4}\\
    b=\boxed{2 \tan \paren{\frac{\pi}{8}}}
\end{gather*}

\subproblem{f}
The volume of the piece with $0\le x\le b$ is larger, as more distance is traveled by the centroid. 
\begin{gather*}
    V= \pi \int_0^\frac{\pi}{8}\cos 2\theta +1\,d\theta\\
    = \pi (\frac{\sin 2\theta}{2}+\theta )\Big|_0^\frac{\pi}{8}\\
    = \boxed{\frac{\pi^2}{8} + \frac{\sqrt{2}\pi}{4}}
\end{gather*}

\problem{11}

\subproblem{a}
\begin{gather*}
    x-y=0\\
    y=x^2 \implies x-x^2=0\\
    d= \boxed{\frac{x-x^2}{\sqrt{2}}}\\
    X-x^2=-(X-x)\\
    X=\frac{x^2+x}{2}\\
    u=\boxed{\frac{x^2+x}{\sqrt{2}}}
\end{gather*}

\subproblem{b}
\begin{gather*}
        V=\int_0^{\sqrt{2}}\pi d^2\, du\\
        =\pi \int_0^1 \paren{\frac{x-x^2}{\sqrt{2}}}^2 \cdot \frac{2x+1}{\sqrt{2}}\, dx\\
        = \frac{\pi}{2\sqrt{2}}\int_0^1 2x^5-3x^4+x^2\, dx\\
        = \frac{\pi}{2\sqrt{2}} ( \frac{2x^6}{6} - \frac{3x^5}{5}+ \frac{x^3}{3})\Big|_0^1\\
        =\boxed{\frac{\pi}{30\sqrt{2}}}
\end{gather*}

\subproblem{c}
\begin{gather*}
    A =\int_0^1 x-x^2\, dx = \frac{x^2}{2}-\frac{x^3}{3}\Big|_0^1=\frac{1}{6}\\
    \delta = \frac{\frac{1}{2}-\frac{2}{5}}{\sqrt{2}} = \frac{1}{10\sqrt{2}}\\
    \frac{\pi}{30\sqrt{2}} = 2\pi \cdot \frac{1}{10\sqrt{2}}\cdot \frac{1}{6}\\
    \boxed{V=2\pi \delta A}
\end{gather*}

\subproblem{d}
\begin{gather*}
    A=1\\
    \delta = \frac{\frac{3}{2}+\frac{1}{2}}{\sqrt{2}} = \sqrt{2}\\
    V=2\pi \cdot \sqrt{2} \cdot 1 = \boxed{2\pi \sqrt{2}}
\end{gather*}
















\end{document}

